% Part: counterfactuals
% Chapter: minimal-change-semantics
% Section: introduction

\documentclass[../../../include/open-logic-section]{subfiles}

\begin{document}

\olfileid[zh]{cnt}{min}{int}

\olsection{引言}

\zhFirst{Stalnaker}{斯达内克}与\zhFirst{Lewis}{刘易斯}提出了关于反事实条件句的理论，例如「如果火柴被划着，它就会亮」。他们的理论是对如何恰当地理解这类句子的真值条件的建议。两个理论背后的想法是这样的：为了判定一个反事实条件句是否为真，我们必须考虑那些为了使前件为真而与现实世界差别最小的可能世界。如果后件在这些可能世界中为真，则反事实条件句为真。例如，假设我手里拿着一根火柴和一个火柴盒。在现实世界中，我只是看着它们，思忖若我划着火柴会发生什么。从现实世界到「我划着火柴」的世界的最小变化，就是我决定去行动并划着火柴的那个世界。说它是最小变化，是因为没有别的改变：我不会同时跳到空中，划火柴也不会点燃我的头发，我不会突然手指乏力，我也不会同时在一场 SuperSoaker 突袭中被水浇个正着，等等。在那个可替代的可能性里，火柴亮了。因此，若我划着火柴，它就会亮，这是真的。

这个直觉性的理论可以与形式化语义配对，用于反事实逻辑。\zhFirst{Lewis}{刘易斯}引入符号「$\boxright$」表示反事实条件句，而\zhFirst{Stalnaker}{斯达内克}使用符号~「$>$」。我们将使用~$\cif$，并把它作为二元联结词添加到命题逻辑中。于是，除了形如 $!A \lif
!B$ 的!!{formula}，我们还拥有形如~$!A \cif !B$ 的!!{formula}。形式化语义与模态逻辑的关系语义一样，建立在这样的模型之上：在其中，!!{formula}在世界处被赋值，而定义 $\mSat{M}{!A \cif !B}[w]$ 的满足条件由某些（另外的）世界~$w'$ 处的 $\mSat{M}{!A}[w']$ 与 $\mSat{M}{!B}[w']$ 给出。哪一些 $w'$？直觉上，是那些满足~$\mSat{M}{!A}[w']$ 且离~$w$ 最近的世界。这要求模型中也得包含一个「接近度」关系。

\zhFirst{Lewis}{刘易斯}引入了一种启发性的方式，以图形方式表示反事实情境。每个可能世界都处于一组包含其他世界的嵌套球面的中心——我们把这些球面画作同心圆。两个球面之间的那些世界离中心世界同等地近，同一个嵌套球面所包含的那些世界更近，而被外围球面包围的那些则更远。
\begin{center}
\begin{tikzpicture}[scale=.7]\small
  \spheresystem{5}
  \draw (0,0) node {$w$};
  \propositionintersect{0}{4}{40}{3.5}
  {\tikzset{proposition/.append={smooth,tension=1.4}}
  \proposition[shift={(1.6,-1)}]{90}{1}{30}{4}}
  \path (1.5,3) node[below] {$\formula{B}$};
  \path (3,0) node {$\formula{A}$};
\end{tikzpicture}
\end{center}
最近的 $!A$-世界是那些满足~$!A$ 且位于中心世界~$w$ 周围最小球面内（灰色区域）的世界 $w'$。直觉上，若在所有最近的 $!A$-世界处 $!B$ 都为真，则 $!A \cif !B$ 在~$w$ 处被满足。

\end{document}
